% =====================================================================
\chapter{A \Mimir{} generálási módszere}
\label{ch:modszer}
% =====================================================================

Ez a fejezet a dolgozat módszertani magja. Először áttekintjük a két összehasonlított csővezetéket --
a naiv RAG-alapvonalat és a tervezett (blueprint) csővezetéket --, majd lépésenként ismertetjük a
darabolás, a tervezés, a visszakeresés, a generálás, az ellenőrzés és a fogalomgráf algoritmusait. A
leírás a kiértékelő keretrendszer implementációját követi
(\code{mimir\_eval/blueprint}), amelynek paraméterei \az{\ref{tab:parameterek}}.~táblázatban szerepelnek.

\section{Áttekintés: naiv RAG és tervezett csővezeték}

\Az{\ref{fig:pipelines}}.~ábra a két csővezetéket veti össze. A \emph{naiv RAG} (\arm{E0}) a
kérésszöveghez legjobban hasonlító három darabot keresi vissza, és egyetlen hívásban kéri a teljes
vizsgát. Ez a rendszer eredeti (v1) működése, és egyben a legolcsóbb ésszerű alapvonal. A \emph{tervezett
csővezeték} (\arm{E1--E3}) ezzel szemben először explicit vizsgatervet készít, majd a terv minden
slotjára külön visszakeresést, egyetlen kérdés generálását és -- az ellenőrző karokban --
ellenőrzést végez.

\begin{figure}[htbp]
\centering
\resizebox{\textwidth}{!}{%
\begin{tikzpicture}[
  node distance=3mm and 3.5mm,
  box/.style={draw, rounded corners=2pt, align=center, font=\footnotesize, minimum height=10mm, inner sep=2pt, text width=20mm},
  hl/.style={box, fill=blue!8, draw=blue!60!black},
  arr/.style={-{Stealth[length=1.6mm]}, semithick},
  lab/.style={font=\footnotesize\bfseries, anchor=west}]
\node[lab] at (-1.3, 1.45) {\arm{E0} -- naiv RAG (alapvonal)};
\node[box] (n1) at (0.4,0.5) {Dokumentum};
\node[box, right=of n1] (n2) {Szemantikus darabok};
\node[box, right=of n2] (n3) {A kéréshez leg\-közelebbi 3 darab};
\node[box, right=of n3, text width=33mm] (n4) {1 LLM-hívás:\\a teljes vizsga};
\node[box, right=of n4] (n5) {Vizsga};
\foreach \a/\b in {n1/n2,n2/n3,n3/n4,n4/n5} \draw[arr] (\a) -- (\b);
\node[lab] at (-1.3, -0.75) {\arm{E1--E3} -- a \Mimir{} tervezett csővezetéke};
\node[box] (b1) at (0.4,-1.75) {Dokumentum és darabok};
\node[hl, right=of b1] (b2) {Fogalomgráf (\arm{E3})};
\node[box, right=of b2] (b3) {Tervező: fogalom, típus, Bloom-szint};
\node[box, right=of b3] (b4) {Slotonkénti keresés (sűrű / hibrid)};
\node[box, right=of b4] (b5) {1 kérdés hivatkozásokkal};
\node[hl, below=6mm of b5] (b6) {Ellenőrző lánc (\arm{E2}, \arm{E3})};
\node[box, right=of b5] (b7) {Összeállítás; oktatói jóváhagyás};
\foreach \a/\b in {b1/b2,b2/b3,b3/b4,b4/b5,b5/b7} \draw[arr] (\a) -- (\b);
\draw[arr] (b5.south) -- (b6.north);
\draw[arr, blue!60!black] (b6.west) -- (b6.west -| b4.south) -- node[left, font=\scriptsize, align=right] {hiba:\\visszajelzés,\\újrapróbálás} (b4.south);
\begin{scope}[on background layer]
\node[draw=gray!60, dashed, rounded corners, fit=(b4)(b5)(b6), inner sep=2.5mm, label={[font=\scriptsize, text=gray]above:minden slotra}] {};
\end{scope}
\end{tikzpicture}}
\caption{A naiv RAG és a \Mimir{} tervezett csővezetéke. Kék: a kísérleti karokban ki- és
bekapcsolt komponensek. A terv a vizsga szerkezetét még az első kérdés megírása előtt rögzíti, és
minden kérdést a saját, hivatkozott kontextusával szemben ellenőriz.}
\label{fig:pipelines}
\end{figure}

A két megközelítés közötti lényegi különbség nem a modell, hanem a feladat felbontása. A naiv
alapvonalban a modellnek egyetlen hívásban kell eldöntenie, hány kérdést, milyen témákról, milyen
nehézségben és milyen disztraktorokkal ír; a tervezett csővezetékben ezek közül a szerkezeti döntéseket
(darabszám, típus, Bloom-szint, téma) determinisztikus kód hozza meg, és a modell feladata egy-egy rövid,
jól körülhatárolt kérdés megírására szűkül, amelyhez csak a releváns kontextust kapja meg.

\begin{table}[htbp]
\centering
\caption{A tervezett csővezeték paraméterei és alapértékei.}
\label{tab:parameterek}
\small
\begin{tabularx}{\textwidth}{@{}>{\raggedright\arraybackslash}p{4.1cm}c>{\raggedright\arraybackslash}X@{}}
\toprule
Paraméter & Alapérték & Jelentés \\
\midrule
\code{retrieval\_k} & 3 & visszakeresett darabok száma slotonként \\
\code{retrieval} & dense & \code{dense}: Bifrost sűrű keresés; \code{hybrid}: BM25 + sűrű, RRF \\
\code{spare\_concepts} & 3 & tartalékfogalmak száma ($n+3$ tervezett fogalom) \\
\code{max\_retries} ($r$) & 1 (pilot: 2) & újrapróbálkozások slotonként visszajelzéssel \\
\code{verifier} & ki/be & ellenőrző lánc (\arm{E1}: ki; \arm{E2}, \arm{E3}: be) \\
\code{blind\_test} & be & vak megoldási próba (\arm{E2f}: ki) \\
\code{leakage\_checks} & be & szivárgás- és kérdésforma-szabályok (a pilot után bevezetve) \\
\code{repair} & be & javító hívás hibás formájú végső kérdésre (a pilot után bevezetve) \\
\code{dedupe\_threshold} & 0,9 & duplikátumküszöb (koszinusz-hasonlóság) \\
\code{overview\_chars} & 12\,000 & a tervezőnek átadott áttekintés karakterkerete \\
\code{max\_context\_chars} & 6\,000 & slotonkénti kontextus karakterkerete \\
\code{graph\_batch\_chars} & 3\,000 & fogalomgráf-kötegek mérete (\arm{E3}) \\
generátor & Qwen2.5-7B & Q4\_K\_M, $T=0{,}2$, 8\,192--16\,384 tokenes kontextus \\
\bottomrule
\end{tabularx}
\end{table}

\section{Szövegkinyerés}

A Wellspring szolgáltatás a PDF-oldalakat először szövegrétegük alapján olvassa ki. Minden oldalra
kiszámítja az olvashatósági arányt,
\begin{equation}
\rho=\frac{|\{\,x\in\text{oldal}:\;x\in\Sigma_{\text{ért}}\,\}|}{|\{\,x\in\text{oldal}:\;x\text{ nem szóköz}\,\}|},
\end{equation}
ahol $\Sigma_{\text{ért}}$ a betűk, számjegyek és szokásos írásjelek halmaza. Ha $\rho<0{,}85$ (hibás
kódolású vagy szkennelt oldal), az oldalt képként dolgozza fel: helyi módban Tesseract OCR-rel,
egyébként egy látás--nyelvi modellel, amely a táblázatokat és listákat is megőrzi. A kiértékelésben
a dokumentumok szövegesek, így ez az ág nem befolyásolja a mérést.

\section{Szemantikus darabolás}
\label{sec:darabolas}

\subsection{Mondatbeágyazás ablakozással}

A darabolás első lépése a mondatfelbontás, amely a mondatvégi írásjelekre és sortörésekre illeszkedő
reguláris kifejezéssel, a karakterpozíciók megőrzésével bontja a szöveget. A mondatbeágyazás a
multilingual-E5-base~\cite{wang2024me5} kódolóval, \emph{kontextuális} módon történik: a modell nem az
egyes mondatokat, hanem hosszabb szövegablakokat kódol, és egy mondat beágyazása a mondatot alkotó
tokenek utolsó rejtett állapotainak átlaga (\emph{mean pooling}). Formálisan, ha a $w$ ablakban a tokenek
rejtett állapotai $h_1,\dots,h_{|w|}$, és $\mathcal{T}(s_i)$ az $s_i$ mondat karakterintervallumával
átfedő tokenek indexhalmaza, akkor
\begin{equation}
\phi(s_i)=\frac{1}{|\mathcal{T}(s_i)|}\sum_{t\in\mathcal{T}(s_i)}h_t .
\end{equation}
Előnye, hogy egy mondat beágyazása a szomszédos mondatok kontextusát is tükrözi, így egy rövid,
önmagában kétértelmű mondat is a témájához közel kerül.

A kódoló legfeljebb 512 tokenes bemenetet fogad. Az eredeti implementáció a teljes dokumentumot
\emph{egyetlen} menetben kódolta, csonkolással; minden mondat, amely a 512. token után kezdődött,
üres tokenhalmazt, így nullvektort kapott. Az új, ablakozott kódoló (\ref{alg:chunk}.~algoritmus) a
mondatokat mondathatáron vágott, legfeljebb 510 tokenes ablakokba csoportosítja, és ablakonként futtatja
a modellt. Rövid, 512 tokennél rövidebb szövegen a két kódoló azonos eredményt ad -- ezt a pilotban
ki is használtuk a futásról futásra jelentkező zaj becslésére (\ref{sec:zaj}.~szakasz).

\begin{algorithm}[htbp]
\caption{Szemantikus darabolás ablakozott mondatbeágyazással}
\label{alg:chunk}
\KwIn{szöveg $\Doc$; módszer $\in\{\text{percentilis},\text{szórás}\}$; küszöbparaméter $\theta$; célméret $\ell^{*}$ (karakter)}
\KwOut{darabok $\Chunks=(c_1,\dots,c_K)$}
$S=(s_1,\dots,s_m)\leftarrow\textsc{Mondatok}(\Doc)$\;
\lIf{$m\le 1$}{\KwRet $(\Doc)$}
\tcp{1. Ablakozás: mondathatáron, legfeljebb 510 token}
$\mathcal{W}\leftarrow\emptyset$; $w\leftarrow()$; $\beta\leftarrow 0$\;
\For{$i=1,\dots,m$}{
  $\ell_i\leftarrow|\textsc{Tokenek}(s_i)|$\;
  \lIf{$w\neq()$ \textbf{és} $\beta+\ell_i>510$}{$\mathcal{W}\leftarrow\mathcal{W}\cup\{w\}$; $w\leftarrow()$; $\beta\leftarrow0$}
  $w\leftarrow w\,\|\,(i)$; $\beta\leftarrow\beta+\ell_i$\;
}
$\mathcal{W}\leftarrow\mathcal{W}\cup\{w\}$\;
\tcp{2. Kontextuális mondatbeágyazás}
\ForEach{$w\in\mathcal{W}$}{
  $H\leftarrow\textsc{Kódoló}(\Doc[a_{w_1}:b_{w_{|w|}}])$ \tcp*{egy menet, $\le 512$ token}
  \ForEach{$i\in w$}{$\phi(s_i)\leftarrow\textsc{Átlag}\{H_t: t\in\mathcal{T}(s_i)\}$ (üres esetén $\mathbf{0}$)}
}
\tcp{3. Küszöb és vágás méretfüggő érzékenységgel}
$\delta_i\leftarrow 1-\cos(\phi(s_i),\phi(s_{i+1}))$ minden $i<m$ esetén\;
$\tau\leftarrow P_{\theta}(\delta)$ ha percentilis, különben $\mu_\delta+\theta\,\sigma_\delta$\;
$c\leftarrow(s_1)$; $\ell\leftarrow|s_1|$; $\Chunks\leftarrow()$\;
\For{$i=1,\dots,m-1$}{
  $\tau_{\text{eff}}\leftarrow\tau/\max(1,\,\ell/\ell^{*})$\;
  \eIf{$\delta_i>\tau_{\text{eff}}$}{$\Chunks\leftarrow\Chunks\,\|\,(c)$; $c\leftarrow(s_{i+1})$; $\ell\leftarrow|s_{i+1}|$}
      {$c\leftarrow c\,\|\,(s_{i+1})$; $\ell\leftarrow\ell+|s_{i+1}|$}
}
\KwRet $\Chunks\,\|\,(c)$\;
\end{algorithm}

\subsection{A vágási szabály}

A vágás két összetevőből áll. Az \emph{alapküszöb} $\tau$ vagy a távolságok $\theta$-adik percentilise
(alapérték: $\theta=85$, vagyis a legnagyobb 15\,\% szomszédos távolságnál vágunk), vagy a
$\mu_\delta+\theta\sigma_\delta$ érték (az \arm{E5b} karban $\theta=1{,}5$). A \emph{méretfüggő
érzékenység} a célméret ($\ell^{*}$) túllépésekor arányosan csökkenti a küszöböt:
\begin{equation}
\tau_{\text{eff}}(\ell)=\frac{\tau}{\max\bigl(1,\,\ell/\ell^{*}\bigr)} .
\end{equation}
Egy célméretnél rövidebb darab tehát csak valódi témaváltásnál zárul, a célméretet $x$-szeresen
meghaladó darab viszont már $\tau/x$ távolságnál is. Ez a „puha korlát” megakadályozza a
túlságosan hosszú darabokat anélkül, hogy a mondatok közepén vágna.

\subsection{A csonkolt kódoló hibájának elemzése}
\label{sec:csonkolt}

A régi kódoló hibája a vágási szabályon keresztül nem nyilvánvaló módon torzít. A 512. token utáni
mondatok nullvektort kapnak; a nullvektor normalizálás után is nulla, így bármely mondattal vett
koszinusz-hasonlósága $0$, a távolság tehát $\delta_i=1$ minden érintett szomszédos párra. Ha a
dokumentum nagyobbik része a határ után van, a távolságok több mint 15\,\%-a $1$, így a percentilis
küszöb $\tau=1$ lesz, és a $\delta_i>\tau$ feltétel a célméret eléréséig soha nem teljesül. A célméret
túllépése után viszont $\tau_{\text{eff}}<1=\delta_i$, így minden további mondatnál vágás történik.
A dokumentum farka ezzel a tartalomtól független, \emph{méretalapú} darabolássá degenerál, a darabok
beágyazása pedig (a mondatbeágyazásoktól függetlenül számolva) már nem tükrözi a vágás logikáját.
$4{,}0$--$4{,}5$ karakter/token aránnyal becsülve a kiértékelő adatkészlet átlagos dokumentumának
68--71\,\%-a esett a határ után, így a hiba a hosszabb dokumentumokon végzett minden korábbi mérést
érintett. Az \arm{E5a} kar a régi kódolót szándékosan megtartja, hogy a hiba hatása a fő vizsgálatban
számszerűsíthető legyen.

\section{Indexelés és sűrű visszakeresés}

A darabokat a Bifrost az E5 kódolóval beágyazza, és egy memóriában futó Qdrant~\cite{qdrant2024}
gyűjteménybe tölti koszinusz-metrikával. A gyűjtemény minden új dokumentum indexelése előtt
törlődik, így a vektortár mindig csak az aktuális feladat tartalmát tartalmazza (\ref{sec:zero}.~szakasz).
A sűrű keresés a $q$ lekérdezésre a
\begin{equation}
\mathrm{top}_k(q)=\argmax_{I\subseteq\{1..K\},\,|I|=k}\;\sum_{i\in I}\cos\bigl(\phi(q),\phi(c_i)\bigr)
\end{equation}
indexhalmazt adja vissza, azaz a $k$ leghasonlóbb darabot.

\section{Az alapvonalak}

\paragraph{Naiv RAG (\arm{E0}).} A lekérdezés a felhasználói kérés szövege (pl.\ „10 feleletválasztós
kérdés, könnyű nehézség”), a kontextus a $\mathrm{top}_3$ darab. A prompt (\ref{app-prompt}.~függelék)
a rendszer éles, v1.0-s promptja: szigorú szabályokat ír elő (zéró hallucináció, a kért darabszám és
típus betartása, egy helyes és három hihető disztraktor, meta-hivatkozás tilalma) és a kimeneti
JSON-sémát. A generálás egyetlen hívás.

\paragraph{Teljes dokumentum (\arm{B-doc-L}, \arm{B-doc-S}).} A status quo modellje: ugyanaz a prompt,
de a kontextus a teljes kinyert dokumentum. A helyi változat 16\,384 tokenes kontextussal és 30\,000
karakterre csonkolt dokumentummal, a szerverváltozat 80\,000 karakterig fut; a csonkolt arányt
($\textit{truncated\_share}$) vizsgánként rögzítjük.

\section{A vizsgaterv}
\label{sec:blueprint}

\subsection{Áttekintés és fogalomtervezés}

A tervező egyetlen hívásban kapja meg a dokumentum \emph{áttekintését}: minden darab azonosítóját
($C_0, C_1,\dots$) és a darab első $\max(120,\lfloor 12\,000/K\rfloor)$ karakterét. Az áttekintés így a
dokumentum hosszától függetlenül korlátos méretű, és minden darab megjelenik benne -- ez a
„lost in the middle” jelenség~\cite{liu2024lost} ellen hat, mert a modell nem egy hosszú, folytonos
szöveget, hanem egy azonosítókkal tagolt, rövid felsorolást lát. A tervező $n+3$ fogalmat ad vissza,
mindegyikhez egy rövid leírást (mit kell kérdezni), a fogalmat alátámasztó darabok azonosítóit és egy
fontossági értéket ($1$: részlet, $2$: fontos, $3$: központi). Ha a tervező válasza értelmezhetetlen,
a csővezeték darabonként egy álfogalomra esik vissza (a darab első nyolc szava), így a terv mindig
létrejön.

\subsection{Bloom-szintek kiosztása}

A nehézséghez rendelt Bloom-eloszlásokat \az{\ref{tab:bloom}}.~táblázat mutatja. A diszkretizálás \az{\ref{all:kvota}}.~állítás szerinti legnagyobb maradék módszerrel történik; a kapott szinteket a
véletlenmaggal megkeverjük (így a szintek nem a fogalmak fontossági sorrendjéhez kötődnek), a kész
vizsgát pedig végül a Bloom-szintek szerint rendezzük, a könnyebbtől a nehezebb felé.

\begin{table}[htbp]
\centering
\caption{A nehézségi szintekhez rendelt Bloom-eloszlások és a $n=10$ kérdéses vizsga kiosztása.}
\label{tab:bloom}
\small
\begin{tabular}{@{}lccccc@{\qquad}c@{}}
\toprule
Nehézség & Emlékezés & Megértés & Alkalmazás & Elemzés & Értékelés & Kiosztás ($n=10$) \\
\midrule
Könnyű   & 0,5 & 0,5 & -- & -- & -- & (5, 5, 0, 0, 0) \\
Közepes  & -- & 0,4 & 0,4 & 0,2 & -- & (0, 4, 4, 2, 0) \\
Nehéz    & -- & -- & 0,3 & 0,4 & 0,3 & (0, 0, 3, 4, 3) \\
\bottomrule
\end{tabular}
\end{table}

\subsection{Mohó, lefedettségvezérelt slotkiosztás}

A slotkiosztás (\ref{alg:blueprint}.~algoritmus) a fogalmakat fontosság szerint rendezi, majd
mohón választ: minden lépésben azt a fogalmat, amely (i) érint még le nem fedett darabot, és ezen belül
(ii) a legfontosabb. \Az{\ref{all:lefed}}.~állítás szerint ez legalább $\min(n,|U|)$ különböző darab
lefedését garantálja. Ha kevesebb fogalom van, mint slot, a fogalmakat körbeforgatva újra
felhasználjuk (eltérő Bloom-szinttel), így a slotok száma mindig pontosan $n$. A ki nem választott
fogalmak alkotják a tartaléklistát.

\begin{algorithm}[htbp]
\caption{Vizsgaterv: Bloom-kiosztás és mohó, lefedettségvezérelt slotkiosztás}
\label{alg:blueprint}
\KwIn{fogalmak $\mathcal{K}$ (mindegyik $D(\kappa)$ darabhalmazzal és $\iota(\kappa)$ fontossággal); $n$; típusok $T$; nehézség $\lambda$; mag $\omega$}
\KwOut{slotok $\mathcal{B}=((\kappa_j,T_j,b_j))_{j=1}^{n}$; tartalékfogalmak $\mathcal{S}$}
\tcp{Bloom-szintek: legnagyobb maradék módszer}
$p\leftarrow\textsc{BloomEloszlás}(\lambda)$; $c_b\leftarrow\lfloor np_b\rfloor$ minden $b$-re\;
\ForEach{$b$ a $np_b-c_b$ törtrész szerint csökkenően, az első $n-\sum_b c_b$ szintre}{$c_b\leftarrow c_b+1$}
$\mathbf{b}\leftarrow\textsc{Kever}_\omega\bigl(\text{minden } b \text{ szint } c_b\text{-szer}\bigr)$\;
\tcp{Mohó lefedés}
$\mathcal{P}\leftarrow\textsc{Rendez}(\mathcal{K},\;-\iota)$; $\mathcal{A}\leftarrow()$; $H\leftarrow\emptyset$\;
\While{$\mathcal{P}\neq\emptyset$ \textbf{és} $|\mathcal{A}|<n$}{
  $\kappa^{+}\leftarrow\argmax_{\kappa\in\mathcal{P}}\bigl(\1[D(\kappa)\setminus H\neq\emptyset],\;\iota(\kappa)\bigr)$ \tcp*{lexikografikus}
  $\mathcal{P}\leftarrow\mathcal{P}\setminus\{\kappa^{+}\}$; $\mathcal{A}\leftarrow\mathcal{A}\,\|\,(\kappa^{+})$; $H\leftarrow H\cup D(\kappa^{+})$\;
}
$i\leftarrow 0$\;
\While(\tcp*[f]{kevés fogalom: újrahasznosítás}){$|\mathcal{A}|<n$}{
  $\mathcal{A}\leftarrow\mathcal{A}\,\|\,(\kappa_{(i \bmod |\mathcal{K}|)})$; $i\leftarrow i+1$\;
}
\KwRet $\mathcal{B}=\bigl((\mathcal{A}_j,\;T_{(j \bmod |T|)},\;\mathbf{b}_j)\bigr)_{j=1}^{n}$; $\mathcal{S}=\mathcal{P}$\;
\end{algorithm}

\section{Slotonkénti visszakeresés}

Minden slot saját kontextust kap. A lekérdezés a fogalom neve és leírása („\emph{fogalom}: \emph{leírás}”),
a kontextus pedig a fogalomhoz a tervező által rendelt első két darab, kiegészítve a keresés
$\mathrm{top}_k$ eredményével (ismétlődés nélkül). A kontextus karakterkerete 6\,000 karakter, hogy a kis
modell „kényelmi zónájában” maradjon: a darabokat sorrendben vesszük fel, amíg a keret engedi.
Ha az ellenőrző azt jelzi, hogy a hivatkozott részlet nem igazolja a kulcsot, a következő kísérletben
$k$ kettővel nő, vagyis a csővezeték szélesebb hálóval tér vissza a visszakereséshez.

\subsection{Hibrid visszakeresés: BM25 és reciprok rangfúzió}

A hibrid kar (\arm{E2h}) a sűrű rangsort egy lexikális BM25-rangsorral~\cite{robertson2009bm25} egyesíti.
A $q$ lekérdezés és a $c$ darab BM25-pontszáma
\begin{equation}
\mathrm{BM25}(q,c)=\sum_{w\in q}\mathrm{idf}(w)\cdot
\frac{f(w,c)\,(k_1+1)}{f(w,c)+k_1\Bigl(1-b+b\,\dfrac{|c|}{\overline{|c|}}\Bigr)},
\qquad
\mathrm{idf}(w)=\ln\Bigl(1+\frac{K-n_w+0{,}5}{n_w+0{,}5}\Bigr),
\end{equation}
ahol $f(w,c)$ a $w$ szó gyakorisága a darabban, $|c|$ a darab tokenhossza, $\overline{|c|}$ az átlagos
darabhossz, $n_w$ a $w$-t tartalmazó darabok száma, $k_1=1{,}5$ és $b=0{,}75$. A magyar nyelv
agglutináló morfológiája miatt a tokenizálás egyszerű, nyelvfüggetlen \emph{prefixtövesítést} alkalmaz:
a hat karakternél hosszabb szavakat hatkarakteres előtagjukra vágja, így az
„immunrendszer”, „immunrendszert”, „immunrendszerben” alakok közös kulcsot kapnak. A két rangsort a
reciprok rangfúzió~\cite{cormack2009rrf} egyesíti:
\begin{equation}
\mathrm{RRF}(c)=\sum_{\rho\in\{\rho_{\text{sűrű}},\,\rho_{\text{BM25}}\}}\frac{1}{k_{\text{rrf}}+\rho(c)},
\qquad k_{\text{rrf}}=60,
\end{equation}
ahol $\rho(c)$ a darab (1-től számozott) rangja az adott rangsorban; a BM25-rangsorból a nulla
pontszámú darabok kimaradnak.

\section{Forráshoz kötött kérdésgenerálás}

A generátor egyetlen kérdést ír a slot terve szerint. A prompt (\ref{app-prompt}.~függelék) tartalmazza
a fogalmat, a kérdezendő tartalmat, a típusszabályt, a Bloom-szint magyar leírását, a nehézséget és a
nyelvet; megismétli az éles prompt szabályait (zéró hallucináció, LaTeX és meta-hivatkozás tilalma);
felsorolja a már elfogadott kérdéseket (legfeljebb az utolsó 12-t), hogy a modell ne ismételje őket;
és megköveteli a \code{citations} mezőt, amelyben a modell megnevezi a kulcsot igazoló darabok
azonosítóit. Ha az előző kísérlet hibás volt, a prompt a hibák listáját is tartalmazza („\emph{Az előző
próbálkozás hibás volt, javítsd:}”). A kérdések egyenkénti generálása kétféle előnnyel jár: a
modellnek egyszerre csak egy döntést kell hoznia, és minden kérdés saját, ellenőrizhető
forrásrészlethez kötődik (G1).

\section{Az ellenőrző lánc}
\label{sec:verifier}

\subsection{Ellenőrzések a költség szerinti sorrendben}

Az ellenőrzések a legolcsóbbtól a legdrágábbig futnak (\ref{tab:checks}.~táblázat); egy olcsó
ellenőrzés hibája esetén a drága, nyelvi modellt igénylő ellenőrzések el sem indulnak. A szabályalapú
ellenőrzések determinisztikusak, a modellalapúak a hivatkozott darabokat (ennek hiányában a teljes
slotkontextust) kapják forrásként.

\begin{table}[htbp]
\centering
\caption{Az ellenőrző lánc lépései végrehajtási sorrendben, a hiba esetén adott visszajelzéssel.}
\label{tab:checks}
\small
\begin{tabularx}{\textwidth}{@{}>{\raggedright\arraybackslash}p{3.4cm}c>{\raggedright\arraybackslash}X@{}}
\toprule
Ellenőrzés & LLM-hívás & Visszajelzés hiba esetén \\
\midrule
Forma ($\mathrm{Shape}$) & -- & feleletválasztós: 4 válasz, 1 helyes \\
Meta-hivatkozás & -- & fogalmazd meg „a szöveg szerint” jellegű hivatkozás nélkül \\
Szivárgás ($L$, $N$) \emph{(új)} & -- & a törzs nem tartalmazhatja a kulcsot, és kérdés legyen \\
Hivatkozás & -- & hivatkozz a kontextus legalább egy darabjára \\
Alátámasztottság és disztraktorok & 1 & a kulcsot a részlet nem igazolja / egy disztraktor igaz lehet / kétértelmű \\
Vak megoldás & 1 & egy független megoldó más választ jelölt meg \\
Duplikátum & -- & túl hasonló egy már elfogadott kérdéshez \\
\bottomrule
\end{tabularx}
\end{table}

Az alátámasztottsági ellenőrzés egyetlen hívásban három ítéletet kér: igazolja-e a forrás a kulcsot
(\code{key\_supported}), hamis-e minden disztraktor (\code{options\_wrong}), és kétértelmű-e a kérdés
(\code{ambiguous}). A vak megoldás külön hívás, amelyben a modell a kulcs ismerete nélkül, véletlen
permutációban látja az opciókat. Ha az ellenőrző modell nem elérhető, a kérdést nem blokkoljuk, de a
kihagyott ellenőrzések számát rögzítjük. Az \arm{E2x} kar az ellenőrző hívásokat egy másik
modellcsaládba tartozó modellel (Llama-3.1-8B~\cite{grattafiori2024llama3}) végzi, hogy elválassza az
önellenőrzés és az ellenőrzés általános hatását.

\subsection{Slotkitöltés visszacsatolással és tartalékokkal}

\Az{\ref{alg:slot}}.~algoritmus egy slot kitöltését írja le. Egy slot legfeljebb $r+1$ kísérletet kap;
minden kísérlet hibalistája a következő kísérlet promptjába kerül. Ha egyik kísérlet sem megy át, a
slot egy tartalékfogalommal újraindul. Ha ez sem sikerül, a \emph{legjobb} kísérlet marad meg,
\emph{nem ellenőrzöttként} megjelölve, hogy az oktató a szerkesztőben lássa. A „legjobb” a
\begin{equation}
\mathrm{rang}(q)=\bigl(\,1-\mathrm{Shape}(q),\;L(q),\;|\text{hibák}(q)|\,\bigr)
\label{eq:rank}
\end{equation}
lexikografikus kulcs szerint legkisebb kísérlet: a formailag helyes kérdés mindig előnyt élvez, ezen belül a
szivárgásmentes, végül a kevesebb hibát tartalmazó. Ha a megmaradt feleletválasztós kérdés formája
mégis hibás, egyetlen \emph{javító} hívás pótolja a hiányzó disztraktorokat a tartalom és a kulcs
megtartásával. \Az{\eqref{eq:rank}} kulcs és a javító hívás a pilot egyik hibájára adott válasz: ott
egy-egy vizsgában hibás formájú végső kérdés maradt meg, ami a vizsga formai megfelelését nullára
állította (\ref{sec:fallback}.~szakasz).

\begin{algorithm}[htbp]
\caption{Egy slot kitöltése ellenőrzéssel, visszacsatolással és tartalékfogalommal}
\label{alg:slot}
\KwIn{slot $(\kappa,T,b)$; elfogadott kérdések $\mathcal{Q}$; tartalékok $\mathcal{S}$; $r$; ellenőrzés be/ki}
\KwOut{kérdés $q$ és \emph{ellenőrzött} jelző}
\SetKwFunction{Fill}{Kitölt}
\SetKwProg{Fn}{Függvény}{:}{}
\Fn{\Fill{$\kappa$}}{
  $\textit{visszajelzés}\leftarrow\varnothing$; $k'\leftarrow 0$; $q^{\text{best}}\leftarrow\bot$\;
  \For{$a=0,\dots,r$}{
    $X\leftarrow\textsc{SlotKontextus}(\kappa,\,k+k')$ \tcp*{$\le$ 6000 karakter}
    $q\leftarrow\textsc{Generál}(\kappa,T,b,X,\mathcal{Q},\textit{visszajelzés})$\;
    \lIf{$q=\bot$}{$\textit{visszajelzés}\leftarrow$ „érvénytelen JSON”; \textbf{folytat}}
    $E\leftarrow\textsc{Ellenőriz}(q,X)$ \tcp*{\ref{tab:checks}.~táblázat, olcsótól drágáig}
    \lIf{$\max_{q'\in\mathcal{Q}}\cos(\phi(q),\phi(q'))\ge0{,}9$}{$E\leftarrow E\cup\{\text{duplikátum}\}$}
    \lIf{$E=\emptyset$}{\KwRet $(q,\textsc{igaz})$}
    \lIf{$\mathrm{rang}(q)<\mathrm{rang}(q^{\text{best}})$}{$q^{\text{best}}\leftarrow q$}
    \lIf{nincs ellenőrzés \textbf{és} $\mathrm{Shape}(q)$ \textbf{és} nem duplikátum}{\KwRet $(q,\textsc{igaz})$}
    $\textit{visszajelzés}\leftarrow E$\;
    \lIf{„nem igazolja” $\in E$}{$k'\leftarrow k'+2$ \tcp*[f]{szélesebb visszakeresés}}
  }
  \KwRet $(q^{\text{best}},\textsc{hamis})$\;
}
$(q,v)\leftarrow$ \Fill{$\kappa$}\;
\If{$\neg v$ \textbf{és} $\mathcal{S}\neq\emptyset$}{
  $(q',v')\leftarrow$ \Fill{$\textsc{Pop}(\mathcal{S})$}\;
  \lIf{$v'$ \textbf{vagy} $q=\bot$ \textbf{vagy} ($\mathrm{Shape}(q')\wedge\neg\mathrm{Shape}(q)$)}{$(q,v)\leftarrow(q',v')$}
}
\If{$T=\textsf{mcq}$ \textbf{és} $\neg\mathrm{Shape}(q)$}{$q\leftarrow\textsc{Javít}(q)$ \tcp*{egy hívás: hiányzó disztraktorok}}
\KwRet $(q,v)$\;
\end{algorithm}

\section{Munkamenet-hatókörű fogalomgráf}
\label{sec:graf}

Az \arm{E3} kar a tervezőt egy fogalomgráffal helyettesíti. A gráfépítés (\ref{alg:graph}.~algoritmus)
a darabokat legfeljebb 3\,000 karakteres kötegekbe csoportosítja, és kötegenként egyetlen hívásban
kinyeri a fogalmakat (definícióval, szülőfogalommal és darabazonosítókkal), az atomi tényeket és a
tipizált relációkat (\emph{is\_a}, \emph{part\_of}, \emph{causes}, \emph{contrasts\_with},
\emph{used\_for}, \emph{related\_to}). A csúcsokat normalizált név alapján, majd a névbeágyazások
$0{,}92$-es koszinusz-küszöbe szerint egyesítjük, ami a magyar ragozott alakokat (pl.\ „T-sejt”, „T-sejtek”)
összevonja. A szülőfogalmak is csúcsok lesznek \emph{is\_a} éllel, így a közös szülő alá tartozó
fogalmak egy lépésben elérhetők egymásból.

Egy $v$ fogalom \emph{centralitása} a fokszám, a darabszám és a tényszám súlyozott összege:
\begin{equation}
\mathrm{cent}(v)=\deg(v)+\tfrac12|D(v)|+\tfrac12|\mathrm{Facts}(v)|,
\end{equation}
amelyből a tervezési fontosság $\iota(v)=3$, ha $\mathrm{cent}(v)\ge4$; $\iota(v)=2$, ha
$\mathrm{cent}(v)\ge2$; egyébként $1$. A gráf három helyen gazdagítja a csővezetéket: (i) a
tervező bemenete a centralitás szerint rendezett fogalomlista; (ii) a disztraktor-ötletek a
\emph{testvérfogalmak}, vagyis azok a $w$ csúcsok, amelyekre $\mathrm{parent}(w)=\mathrm{parent}(v)$
vagy $N(w)\cap N(v)\ne\emptyset$, de $w\notin N(v)$ és $w$ nem $v$ szülője -- ezek ugyanabba a
kategóriába tartoznak, ezért hihetők, de nem azonosak a kulccsal; (iii) a kontextus kiegészül a
szomszédos fogalmak legtöbbször előforduló, saját darabhalmazon kívüli két darabjával (1-ugrásos
szomszédság).

\begin{algorithm}[htbp]
\caption{Munkamenet-hatókörű fogalomgráf építése}
\label{alg:graph}
\KwIn{darabok $\Chunks$; kötegméret $B=3000$; hasonlósági küszöb $\eta=0{,}92$}
\KwOut{gráf $\mathcal{G}=(V,E)$ fogalmakkal, tényekkel és darabhalmazokkal}
$V\leftarrow\emptyset$; $E\leftarrow\emptyset$\;
\ForEach{$\mathcal{I}\in\textsc{Kötegek}(\Chunks,B)$}{
  $(\textit{Fog},\textit{Tény},\textit{Rel})\leftarrow\textsc{LLM}_{\text{gráf}}\bigl(\{[C_i]\,c_i: i\in\mathcal{I}\}\bigr)$ \tcp*{1 hívás}
  \ForEach{$(\textit{név},\textit{def},\textit{szülő},D)\in\textit{Fog}$}{$\textsc{Hozzáad}(V,\textsc{Norm}(\textit{név}),\textit{def},\textit{szülő},D)$ \tcp*{egyesít, ha létezik}}
  \ForEach{$(\textit{szöveg},\textit{fogalmak},i)\in\textit{Tény}$}{minden érintett fogalomhoz csatolja a tényt és az $i$ darabot}
  \ForEach{$(u,\textit{típus},w)\in\textit{Rel}$}{$E\leftarrow E\cup\{\{u,w\}\}$}
  \ForEach{$v\in V$ szülővel}{$E\leftarrow E\cup\{\{v,\mathrm{parent}(v)\}\}$ \tcp*{\emph{is\_a} él}}
}
\ForEach{$u,w\in V$, $u\prec w$}{
  \lIf{$\cos(\phi(\nu_u),\phi(\nu_w))\ge\eta$}{$\textsc{Összevon}(u,w)$ \tcp*{$\nu_u$: $u$ neve; darabok, tények, élek uniója}}
}
\KwRet $\mathcal{G}$\;
\end{algorithm}

A gráf csak a vizsga élettartamáig létezik a memóriában, összhangban a zéró megőrzéssel. Költsége
$\lceil|\Doc|/B\rceil$ nagyságrendű további hívás; egy 1\,700 karakteres dokumentumon ez egyetlen hívás
volt, amely a pilot három futásában 14, 18 és 31 fogalmat, illetve 3, 9 és 19 relációt adott -- a
szórás maga is jelzi, hogy a kis modell gráfkinyerése instabil.

\section{Összeállítás és emberi felügyelet}

Az elfogadott kérdéseket a Bloom-szint szerint rendezzük, és minden kérdéshez rögzítjük a hivatkozott
darabokat, a kontextus darabjait, a tervezett fogalmat és az \emph{ellenőrzött} jelzőt. A vázlat egy
szerkesztőbe kerül, ahol az oktató szerkesztheti, átrendezheti és törölheti a kérdéseket, módosíthatja
a kulcsot, egyetlen kérdést újragenerálhat, és minden kérdés mellett megtekintheti a hivatkozott
forrásrészletet. Jóváhagyás nélkül semmi sem exportálható, az exportált dokumentum pedig feltünteti,
hogy a tartalom MI-rendszer közreműködésével készült, és melyik modellel (G5).
